\documentclass[11pt,a4paper]{article}

% ─── Packages ─────────────────────────────────────────────────────────────────
\usepackage[T1]{fontenc}
\usepackage[utf8]{inputenc}
\usepackage[margin=2.4cm]{geometry}
\usepackage{microtype}
\usepackage{lmodern}
\usepackage{xcolor}
\usepackage{listings}
\usepackage{booktabs}
\usepackage{tabularx}
\usepackage{longtable}
\usepackage{hyperref}
\usepackage{titlesec}
\usepackage{parskip}
\usepackage{fancyhdr}
\usepackage{mdframed}
\usepackage{enumitem}
\usepackage{amsmath}

% ─── Colours ──────────────────────────────────────────────────────────────────
\definecolor{xoop-blue}    {HTML}{1565C0}
\definecolor{xoop-green}   {HTML}{2E7D32}
\definecolor{xoop-red}     {HTML}{B71C1C}
\definecolor{xoop-gray}    {HTML}{546E7A}
\definecolor{xoop-bg}      {HTML}{F5F7FA}
\definecolor{xoop-border}  {HTML}{CFD8DC}
\definecolor{xoop-comment} {HTML}{78909C}
\definecolor{xoop-string}  {HTML}{388E3C}
\definecolor{xoop-keyword} {HTML}{1565C0}
\definecolor{xoop-attr}    {HTML}{6A1B9A}

% ─── Listings style ───────────────────────────────────────────────────────────
\lstdefinelanguage{xoop}{
  morecomment=[s]{<!--}{-->},
  morestring=[b]",
  sensitive=false,
  morekeywords={
    XoopProgram, Using, Namespace, Class, Interface, Enum,
    Field, Property, Constructor, Method, Parameter, Get, Set,
    Body, Member, Base, This, Arg
  },
}

\lstdefinestyle{xoopstyle}{
  language=xoop,
  backgroundcolor=\color{xoop-bg},
  basicstyle=\ttfamily\small,
  keywordstyle=\color{xoop-keyword}\bfseries,
  stringstyle=\color{xoop-string},
  commentstyle=\color{xoop-comment}\itshape,
  identifierstyle=\color{xoop-attr},
  breaklines=true,
  breakatwhitespace=true,
  showstringspaces=false,
  tabsize=2,
  frame=single,
  rulecolor=\color{xoop-border},
  framesep=4pt,
  xleftmargin=8pt,
  xrightmargin=4pt,
  numbers=left,
  numberstyle=\tiny\color{xoop-gray},
  numbersep=8pt,
  captionpos=b,
}

\lstdefinestyle{csstyle}{
  language={[Sharp]C},
  backgroundcolor=\color{xoop-bg},
  basicstyle=\ttfamily\small,
  keywordstyle=\color{xoop-keyword}\bfseries,
  stringstyle=\color{xoop-string},
  commentstyle=\color{xoop-comment}\itshape,
  breaklines=true,
  showstringspaces=false,
  tabsize=4,
  frame=single,
  rulecolor=\color{xoop-border},
  framesep=4pt,
  xleftmargin=8pt,
  xrightmargin=4pt,
  numbers=left,
  numberstyle=\tiny\color{xoop-gray},
  numbersep=8pt,
  captionpos=b,
}

\lstset{style=xoopstyle}

% ─── Section formatting ───────────────────────────────────────────────────────
\titleformat{\section}{\Large\bfseries\color{xoop-blue}}{}{0em}{}[\titlerule]
\titleformat{\subsection}{\large\bfseries\color{xoop-gray}}{}{0em}{}
\titleformat{\subsubsection}{\normalsize\bfseries}{}{0em}{}

% ─── Header / footer ─────────────────────────────────────────────────────────
\pagestyle{fancy}
\fancyhf{}
\renewcommand{\headrulewidth}{0.4pt}
\fancyhead[L]{\small\textcolor{xoop-gray}{XOOP Language Reference}}
\fancyhead[R]{\small\textcolor{xoop-gray}{v1.0.0}}
\fancyfoot[C]{\small\thepage}

% ─── Hyperref ─────────────────────────────────────────────────────────────────
\hypersetup{
  colorlinks=true,
  linkcolor=xoop-blue,
  urlcolor=xoop-blue,
  pdftitle={XOOP Language Reference},
  pdfauthor={XOOP Compiler},
}

% ─── Convenience macros ───────────────────────────────────────────────────────
\newcommand{\el}[1]{\texttt{\textcolor{xoop-keyword}{<#1>}}}
\newcommand{\attr}[1]{\texttt{\textcolor{xoop-attr}{#1}}}
\newcommand{\val}[1]{\texttt{\textcolor{xoop-string}{"#1"}}}
\newcommand{\cs}[1]{\texttt{#1}}
\newcommand{\required}{\textcolor{xoop-red}{\textbf{required}}}
\newcommand{\optional}{\textcolor{xoop-green}{optional}}
\newcommand{\default}[1]{\textcolor{xoop-green}{default: \texttt{#1}}}

\newenvironment{note}{%
  \begin{mdframed}[backgroundcolor=xoop-bg,linecolor=xoop-border,
                   innertopmargin=6pt,innerbottommargin=6pt,
                   innerleftmargin=10pt,innerrightmargin=6pt]
  \small\textbf{Note:}\ }{%
  \end{mdframed}}

% ═════════════════════════════════════════════════════════════════════════════
\begin{document}

% ─── Title page ───────────────────────────────────────────────────────────────
\begin{titlepage}
  \centering
  \vspace*{3cm}
  {\fontsize{48}{56}\selectfont\bfseries\color{xoop-blue} XOOP \par}
  \vspace{0.5em}
  {\Large\color{xoop-gray} XML Object-Oriented Programming Language \par}
  \vspace{2em}
  \rule{0.5\textwidth}{0.8pt} \par
  \vspace{1.5em}
  {\large Language Reference \& Compiler Guide \par}
  \vspace{0.5em}
  {\normalsize Version 1.0.0 \quad\textbullet\quad April 2026 \par}
  \vfill
  {\small Compiled to C\# via the XOOP Compiler. \par}
  {\small Source: \texttt{.xoop} \quad\textrightarrow\quad \texttt{.cs}}
\end{titlepage}

\tableofcontents
\newpage

% ═════════════════════════════════════════════════════════════════════════════
\section{Introduction}

\textbf{XOOP} (\emph{XML Object-Oriented Programming}) is a source-to-source
compiled language that uses well-formed XML as its surface syntax.
A \texttt{.xoop} file is valid XML; the XOOP compiler parses it into an abstract
syntax tree (AST) and emits idiomatic C\# source code.

The core design goals are:

\begin{enumerate}[nosep]
  \item \textbf{Toolability} — XML is natively understood by editors, validators,
        linters, and diff tools with no additional grammar files.
  \item \textbf{Verbosity as clarity} — every declaration is explicit, making
        codebases easy to audit programmatically (XSLT, XPath, etc.).
  \item \textbf{Full OOP parity} — classes, interfaces, enums, inheritance,
        generics, access modifiers, and async all map cleanly to C\# constructs.
  \item \textbf{Zero runtime} — the compiler output is plain C\# with no XOOP
        dependency; compiled programs depend only on the .NET runtime.
\end{enumerate}

% ─── Architecture overview ────────────────────────────────────────────────────
\subsection{Compiler Architecture}

The compiler is composed of three layers:

\begin{center}
\begin{tabular}{lll}
\toprule
\textbf{Layer} & \textbf{Class} & \textbf{Responsibility} \\
\midrule
Parser  & \cs{XoopParser}    & XML $\rightarrow$ AST (\cs{ProgramNode} tree) \\
Emitter & \cs{CSharpEmitter} & AST $\rightarrow$ formatted C\# source \\
Driver  & \cs{XoopCompiler}  & Orchestrates parser $+$ emitter \\
\bottomrule
\end{tabular}
\end{center}

\begin{note}
Method and constructor bodies are written in plain C\# inside
\texttt{<Body>} elements. CDATA sections (\texttt{<![CDATA[\ \ldots\ ]]>}) are
strongly recommended to avoid XML escaping issues inside bodies.
\end{note}

% ═════════════════════════════════════════════════════════════════════════════
\section{File Structure}

Every XOOP source file must be a valid XML document whose root element is
\el{XoopProgram}.

\begin{lstlisting}[caption={Minimal XOOP file skeleton}]
<?xml version="1.0" encoding="utf-8"?>
<XoopProgram>

  <!-- namespace imports -->
  <Using namespace="System"/>

  <!-- optional namespace block -->
  <Namespace name="MyApp">
    <Class name="Hello">
      <!-- members... -->
    </Class>
  </Namespace>

</XoopProgram>
\end{lstlisting}

Top-level children of \el{XoopProgram} (in any order):

\begin{itemize}[nosep]
  \item \el{Using} — namespace import
  \item \el{Namespace} — logical grouping
  \item \el{Class} — top-level class (no enclosing namespace)
  \item \el{Interface} — top-level interface
  \item \el{Enum} — top-level enum
\end{itemize}

% ═════════════════════════════════════════════════════════════════════════════
\section{Element Reference}

% ─────────────────────────────────────────────────────────────────────────────
\subsection{\el{Using} — Namespace Import}

Emits a \cs{using} directive.

\begin{center}
\begin{tabular}{lll}
\toprule
Attribute & Req? & Description \\
\midrule
\attr{namespace} (or \attr{ns}) & \required & The namespace to import \\
\bottomrule
\end{tabular}
\end{center}

\begin{lstlisting}[caption={Importing namespaces}]
<Using namespace="System"/>
<Using namespace="System.Collections.Generic"/>
<Using ns="System.Linq"/>
\end{lstlisting}

Emits:
\begin{lstlisting}[style=csstyle, numbers=none]
using System;
using System.Collections.Generic;
using System.Linq;
\end{lstlisting}

% ─────────────────────────────────────────────────────────────────────────────
\subsection{\el{Namespace}}

Wraps declarations in a C\# \cs{namespace} block.

\begin{center}
\begin{tabular}{lll}
\toprule
Attribute & Req? & Description \\
\midrule
\attr{name} & \required & Fully qualified namespace name \\
\bottomrule
\end{tabular}
\end{center}

Children: \el{Class}, \el{Interface}, \el{Enum}.

\begin{lstlisting}[caption={Namespace declaration}]
<Namespace name="MyApp.Services">
  <Class name="UserService"> ... </Class>
</Namespace>
\end{lstlisting}

% ─────────────────────────────────────────────────────────────────────────────
\subsection{\el{Class}}

Declares a class. Nested \el{Class} and \el{Enum} elements produce nested types.

\begin{center}
\begin{tabularx}{\linewidth}{llX}
\toprule
Attribute & Req? & Description \\
\midrule
\attr{name}       & \required & Class name \\
\attr{access}     & \optional & \val{public} \textbar{} \val{private} \textbar{} \val{internal} \textbar{} \val{protected} \quad \default{public} \\
\attr{extends}    & \optional & Single base class name \\
\attr{implements} & \optional & Comma-separated list of interfaces \\
\attr{generics}   & \optional & Comma-separated generic type parameters (e.g.\ \val{T,U}) \\
\attr{abstract}   & \optional & \val{true} / \val{false} \quad \default{false} \\
\attr{sealed}     & \optional & \val{true} / \val{false} \quad \default{false} \\
\attr{static}     & \optional & \val{true} / \val{false} \quad \default{false} \\
\attr{partial}    & \optional & \val{true} / \val{false} \quad \default{false} \\
\bottomrule
\end{tabularx}
\end{center}

Children (in recommended declaration order):
\el{Field}, \el{Property}, \el{Constructor}, \el{Method}, \el{Class}, \el{Enum}.

\begin{lstlisting}[caption={Class with inheritance, interface, and generics}]
<Class name="Repository"
       access="public"
       extends="BaseRepository"
       implements="IRepository,IDisposable"
       generics="T">

  <Field name="_context" type="DbContext" access="private" readonly="true"/>
  <Constructor access="public">
    <Parameter name="context" type="DbContext"/>
    <Base><Arg>context</Arg></Base>
  </Constructor>
  <Method name="FindById" returnType="T?" access="public">
    <Parameter name="id" type="int"/>
    <Body><![CDATA[
      return _context.Set<T>().Find(id);
    ]]></Body>
  </Method>

</Class>
\end{lstlisting}

% ─────────────────────────────────────────────────────────────────────────────
\subsection{\el{Interface}}

Declares an interface containing method signatures and property signatures.

\begin{center}
\begin{tabularx}{\linewidth}{llX}
\toprule
Attribute & Req? & Description \\
\midrule
\attr{name}    & \required & Interface name \\
\attr{access}  & \optional & \default{public} \\
\attr{extends} & \optional & Comma-separated parent interface names \\
\bottomrule
\end{tabularx}
\end{center}

Children: \el{Method} (signature only, no \el{Body}), \el{Property}
(with \attr{get} / \attr{set} boolean attributes).

\begin{lstlisting}[caption={Interface declaration}]
<Interface name="IAnimal" extends="IComparable,IDisposable">
  <Property name="Name" type="string" get="true" set="false"/>
  <Method name="Speak"    returnType="void"/>
  <Method name="GetSound" returnType="string">
    <Parameter name="loud" type="bool" default="false"/>
  </Method>
</Interface>
\end{lstlisting}

Emits:
\begin{lstlisting}[style=csstyle, numbers=none]
public interface IAnimal : IComparable, IDisposable
{
    string Name { get; }
    void Speak();
    string GetSound(bool loud = false);
}
\end{lstlisting}

% ─────────────────────────────────────────────────────────────────────────────
\subsection{\el{Enum}}

Declares an enumeration.

\begin{center}
\begin{tabularx}{\linewidth}{llX}
\toprule
Attribute & Req? & Description \\
\midrule
\attr{name}   & \required & Enum name \\
\attr{access} & \optional & \default{public} \\
\attr{type}   & \optional & Underlying numeric type (e.g.\ \val{byte}, \val{int}) \\
\bottomrule
\end{tabularx}
\end{center}

Each member is an \el{Member} element with:

\begin{center}
\begin{tabular}{lll}
\toprule
Attribute & Req? & Description \\
\midrule
\attr{name}  & \required & Member name \\
\attr{value} & \optional & Explicit integer value \\
\bottomrule
\end{tabular}
\end{center}

\begin{lstlisting}[caption={Enum with explicit values}]
<Enum name="HttpStatus" type="int">
  <Member name="OK"                  value="200"/>
  <Member name="BadRequest"          value="400"/>
  <Member name="Unauthorized"        value="401"/>
  <Member name="NotFound"            value="404"/>
  <Member name="InternalServerError" value="500"/>
</Enum>
\end{lstlisting}

% ─────────────────────────────────────────────────────────────────────────────
\subsection{\el{Field}}

Declares a class-level field.

\begin{center}
\begin{tabularx}{\linewidth}{llX}
\toprule
Attribute & Req? & Description \\
\midrule
\attr{name}     & \required & Field name \\
\attr{type}     & \required & C\# type \\
\attr{access}   & \optional & \default{private} \\
\attr{static}   & \optional & \default{false} \\
\attr{readonly} & \optional & \default{false} \\
\attr{const}    & \optional & \default{false} — implies \texttt{static} semantics \\
\attr{default}  & \optional & Initialiser expression (C\# literal or expression) \\
\bottomrule
\end{tabularx}
\end{center}

\begin{lstlisting}[caption={Field examples}]
<Field name="_id"      type="int"    access="private" readonly="true"/>
<Field name="_name"    type="string" access="private" readonly="true" default='"unknown"'/>
<Field name="MaxItems" type="int"    access="public"  const="true" default="100"/>
<Field name="_count"   type="int"    access="protected" static="true" default="0"/>
\end{lstlisting}

Emits:
\begin{lstlisting}[style=csstyle, numbers=none]
private readonly int _id;
private readonly string _name = "unknown";
public const int MaxItems = 100;
protected static int _count = 0;
\end{lstlisting}

% ─────────────────────────────────────────────────────────────────────────────
\subsection{\el{Property}}

XOOP supports three property forms:

\subsubsection{Auto-property (attribute-only)}

\begin{lstlisting}[caption={Auto-property}]
<Property name="Count" type="int" access="public"
          setAccess="private" default="0"/>
\end{lstlisting}

Emits:
\begin{lstlisting}[style=csstyle, numbers=none]
public int Count { get; private set; } = 0;
\end{lstlisting}

\subsubsection{Field-backed property (\attr{fieldRef})}

\begin{lstlisting}[caption={Field-backed read-only property}]
<Field name="_name" type="string" access="private" readonly="true"/>
<Property name="Name" type="string" access="public"
          fieldRef="_name" set="false"/>
\end{lstlisting}

Emits:
\begin{lstlisting}[style=csstyle, numbers=none]
public string Name
{
    get => _name;
}
\end{lstlisting}

\subsubsection{Custom body (\el{Get} / \el{Set} children)}

\begin{lstlisting}[caption={Property with custom accessors}]
<Property name="DisplayName" type="string" access="public">
  <Get><![CDATA[ return _first + " " + _last; ]]></Get>
  <Set><![CDATA[
    var parts = value.Split(' ', 2);
    _first = parts[0];
    _last  = parts.Length > 1 ? parts[1] : "";
  ]]></Set>
</Property>
\end{lstlisting}

\begin{center}
\begin{tabularx}{\linewidth}{llX}
\toprule
Attribute & Req? & Description \\
\midrule
\attr{name}      & \required & Property name \\
\attr{type}      & \required & C\# type \\
\attr{access}    & \optional & \default{public} \\
\attr{setAccess} & \optional & Narrower access for setter (e.g.\ \val{private}) \\
\attr{get}       & \optional & Include getter? \default{true} \\
\attr{set}       & \optional & Include setter? \default{true} \\
\attr{fieldRef}  & \optional & Backing field — generates expression-body accessors \\
\attr{default}   & \optional & Auto-property initialiser \\
\attr{static}    & \optional & \default{false} \\
\attr{virtual}   & \optional & \default{false} \\
\attr{override}  & \optional & \default{false} \\
\bottomrule
\end{tabularx}
\end{center}

% ─────────────────────────────────────────────────────────────────────────────
\subsection{\el{Constructor}}

\begin{center}
\begin{tabularx}{\linewidth}{llX}
\toprule
Attribute & Req? & Description \\
\midrule
\attr{access} & \optional & \default{public} \\
\bottomrule
\end{tabularx}
\end{center}

Children:
\begin{itemize}[nosep]
  \item \el{Parameter} (zero or more)
  \item \el{Base} — emits \cs{: base(\ldots)}; contains \el{Arg} children
  \item \el{This} — emits \cs{: this(\ldots)}; contains \el{Arg} children
  \item \el{Body} — C\# method body
\end{itemize}

\begin{lstlisting}[caption={Constructor with base call}]
<Class name="PremiumUser" extends="User">
  <Constructor access="public">
    <Parameter name="id"   type="int"/>
    <Parameter name="name" type="string"/>
    <Base><Arg>id</Arg><Arg>name</Arg></Base>
    <Body><![CDATA[
      IsPremium = true;
    ]]></Body>
  </Constructor>
</Class>
\end{lstlisting}

Emits:
\begin{lstlisting}[style=csstyle, numbers=none]
public PremiumUser(int id, string name) : base(id, name)
{
    IsPremium = true;
}
\end{lstlisting}

% ─────────────────────────────────────────────────────────────────────────────
\subsection{\el{Method}}

\begin{center}
\begin{tabularx}{\linewidth}{llX}
\toprule
Attribute & Req? & Description \\
\midrule
\attr{name}       & \required & Method name \\
\attr{returnType} & \optional & \default{void} \\
\attr{access}     & \optional & \default{public} \\
\attr{generics}   & \optional & Comma-separated type parameters \\
\attr{static}     & \optional & \default{false} \\
\attr{virtual}    & \optional & \default{false} \\
\attr{override}   & \optional & \default{false} \\
\attr{abstract}   & \optional & \default{false} — omits body braces, emits \texttt{;} \\
\attr{sealed}     & \optional & \default{false} \\
\attr{new}        & \optional & \default{false} — emits \cs{new} modifier \\
\attr{async}      & \optional & \default{false} — emits \cs{async} modifier \\
\bottomrule
\end{tabularx}
\end{center}

Children: \el{Parameter} (zero or more), \el{Body}.

\begin{lstlisting}[caption={Method examples}]
<!-- Abstract method (no body) -->
<Method name="GetArea" returnType="double" abstract="true"/>

<!-- Virtual method -->
<Method name="ToString" returnType="string" override="true">
  <Body><![CDATA[ return $"Point({X}, {Y})"; ]]></Body>
</Method>

<!-- Generic static method -->
<Method name="Create" returnType="T" static="true" generics="T"
        access="public">
  <Parameter name="args" type="object[]" params="true"/>
  <Body><![CDATA[
    return (T)Activator.CreateInstance(typeof(T), args)!;
  ]]></Body>
</Method>

<!-- Async method -->
<Method name="FetchAsync" returnType="Task&lt;string&gt;" async="true">
  <Parameter name="url" type="string"/>
  <Body><![CDATA[
    using var client = new HttpClient();
    return await client.GetStringAsync(url);
  ]]></Body>
</Method>
\end{lstlisting}

% ─────────────────────────────────────────────────────────────────────────────
\subsection{\el{Parameter}}

\begin{center}
\begin{tabularx}{\linewidth}{llX}
\toprule
Attribute & Req? & Description \\
\midrule
\attr{name}    & \required & Parameter name \\
\attr{type}    & \required & C\# type \\
\attr{default} & \optional & Default value expression \\
\attr{params}  & \optional & \val{true} — prepends \cs{params} keyword \\
\attr{ref}     & \optional & \val{true} — prepends \cs{ref} keyword \\
\attr{out}     & \optional & \val{true} — prepends \cs{out} keyword \\
\attr{in}      & \optional & \val{true} — prepends \cs{in} keyword \\
\bottomrule
\end{tabularx}
\end{center}

\begin{lstlisting}[caption={Parameter modifiers}]
<Parameter name="count"  type="int"      default="0"/>
<Parameter name="values" type="int[]"    params="true"/>
<Parameter name="result" type="string"   out="true"/>
<Parameter name="data"   type="byte[]"   ref="true"/>
\end{lstlisting}

% ═════════════════════════════════════════════════════════════════════════════
\section{Access Modifiers}

All \attr{access} attributes accept the following values:

\begin{center}
\begin{tabular}{ll}
\toprule
Value & C\# keyword \\
\midrule
\val{public}            & \cs{public} \\
\val{private}           & \cs{private} \\
\val{protected}         & \cs{protected} \\
\val{internal}          & \cs{internal} \\
\val{protected internal}& \cs{protected internal} \\
\val{private protected} & \cs{private protected} \\
\bottomrule
\end{tabular}
\end{center}

The default access for class members follows the same convention as C\#:
fields default to \cs{private}, methods and properties default to \cs{public}.

% ═════════════════════════════════════════════════════════════════════════════
\section{Boolean Attributes}

Any boolean attribute (\attr{abstract}, \attr{static}, \attr{readonly}, etc.)
accepts the following truthy/falsy values (case-insensitive):

\begin{center}
\begin{tabular}{ll}
\toprule
Truthy & Falsy \\
\midrule
\val{true}, \val{yes}, \val{1} & \val{false}, \val{no}, \val{0} \\
\bottomrule
\end{tabular}
\end{center}

Omitting a boolean attribute is equivalent to \val{false}.

% ═════════════════════════════════════════════════════════════════════════════
\section{Method Bodies}

Method, constructor, getter, and setter bodies are written in plain C\# inside
the \el{Body}, \el{Get}, or \el{Set} elements respectively.

\subsection{CDATA Recommendation}

Wrap body content in a CDATA section to avoid XML escaping interference:

\begin{lstlisting}[caption={Using CDATA for method bodies}]
<Method name="Print" returnType="void">
  <Parameter name="value" type="int"/>
  <Body><![CDATA[
    if (value > 0)
        Console.WriteLine($"Positive: {value}");
    else if (value < 0)
        Console.WriteLine($"Negative: {value}");
    else
        Console.WriteLine("Zero");
  ]]></Body>
</Method>
\end{lstlisting}

\subsection{Without CDATA (XML Escaping Required)}

If CDATA is not used, XML special characters must be escaped:

\begin{center}
\begin{tabular}{ll}
\toprule
Character & XML escape \\
\midrule
\texttt{<}  & \texttt{\&lt;} \\
\texttt{>}  & \texttt{\&gt;} \\
\texttt{\&} & \texttt{\&amp;} \\
\texttt{"}  & \texttt{\&quot;} \\
\texttt{'}  & \texttt{\&apos;} \\
\bottomrule
\end{tabular}
\end{center}

The emitter automatically strips common leading whitespace (the CDATA indentation
offset), so bodies can be indented naturally without affecting the emitted code.

% ═════════════════════════════════════════════════════════════════════════════
\section{Complete Example — Hello World}

\begin{lstlisting}[caption={examples/hello.xoop}]
<?xml version="1.0" encoding="utf-8"?>
<XoopProgram>

  <Using namespace="System"/>

  <Namespace name="Hello">

    <Class name="Greeter" access="public">

      <Field name="_greeting" type="string" access="private"
             readonly="true" default='"Hello"'/>

      <Constructor access="public">
        <Parameter name="greeting" type="string" default='"Hello"'/>
        <Body><![CDATA[
          _greeting = greeting;
        ]]></Body>
      </Constructor>

      <Method name="Greet" returnType="void" access="public">
        <Parameter name="name" type="string"/>
        <Body><![CDATA[
          Console.WriteLine($"{_greeting}, {name}!");
        ]]></Body>
      </Method>

    </Class>

    <Class name="Program" access="public" static="true">
      <Method name="Main" returnType="void" access="public" static="true">
        <Parameter name="args" type="string[]"/>
        <Body><![CDATA[
          var greeter = new Greeter();
          greeter.Greet("World");

          var custom = new Greeter("Howdy");
          custom.Greet("partner");
        ]]></Body>
      </Method>
    </Class>

  </Namespace>

</XoopProgram>
\end{lstlisting}

Generated C\# output:
\begin{lstlisting}[style=csstyle, caption={Generated hello.cs (excerpt)}]
using System;

namespace Hello
{
    public class Greeter
    {
        private readonly string _greeting = "Hello";

        public Greeter(string greeting = "Hello")
        {
            _greeting = greeting;
        }

        public void Greet(string name)
        {
            Console.WriteLine($"{_greeting}, {name}!");
        }
    }

    public static class Program
    {
        public static void Main(string[] args)
        {
            var greeter = new Greeter();
            greeter.Greet("World");

            var custom = new Greeter("Howdy");
            custom.Greet("partner");
        }
    }
}
\end{lstlisting}

% ═════════════════════════════════════════════════════════════════════════════
\section{Complete Example — Animal Kingdom (Full OOP Demo)}

This example demonstrates the full OOP feature set: abstract classes,
interfaces, enums, inheritance chains, field-backed read-only properties,
virtual/override methods, and constructor chaining.

\begin{lstlisting}[caption={examples/animals.xoop (abridged)}]
<XoopProgram>
  <Using namespace="System"/>
  <Using namespace="System.Collections.Generic"/>

  <Namespace name="AnimalKingdom">

    <Interface name="ISpeakable">
      <Method name="Speak"    returnType="void"/>
      <Method name="GetSound" returnType="string"/>
    </Interface>

    <Enum name="AnimalKind">
      <Member name="Unknown" value="0"/>
      <Member name="Dog"/>
      <Member name="Cat"/>
      <Member name="Bird"/>
    </Enum>

    <Class name="Animal" abstract="true" implements="ISpeakable">
      <Field name="_name"  type="string" access="private"   readonly="true"/>
      <Field name="_sound" type="string" access="protected"/>

      <Constructor access="public">
        <Parameter name="name"  type="string"/>
        <Parameter name="sound" type="string"/>
        <Body><![CDATA[ _name = name; _sound = sound; ]]></Body>
      </Constructor>

      <Property name="Name" type="string" fieldRef="_name" set="false"/>
      <Method name="Kind"   returnType="AnimalKind" abstract="true"/>
      <Method name="Speak"  returnType="void" virtual="true">
        <Body><![CDATA[
          Console.WriteLine($"{_name} says: {_sound}!");
        ]]></Body>
      </Method>
    </Class>

    <Class name="Dog" extends="Animal">
      <Constructor access="public">
        <Parameter name="name" type="string"/>
        <Base><Arg>name</Arg><Arg>"Woof"</Arg></Base>
      </Constructor>
      <Method name="Kind" returnType="AnimalKind" override="true">
        <Body><![CDATA[ return AnimalKind.Dog; ]]></Body>
      </Method>
    </Class>

    <!-- Cat and Bird follow the same pattern... -->

  </Namespace>
</XoopProgram>
\end{lstlisting}

% ═════════════════════════════════════════════════════════════════════════════
\section{Using the Compiler}

\subsection{Build}

\begin{verbatim}
dotnet build
\end{verbatim}

\subsection{Compile a \texttt{.xoop} file}

\begin{verbatim}
dotnet run -- <input.xoop> [output.cs]
\end{verbatim}

If \texttt{output.cs} is omitted, the output file is placed next to the input
file with the extension changed to \texttt{.cs}.

\subsection{Examples}

\begin{verbatim}
dotnet run -- examples/hello.xoop
dotnet run -- examples/animals.xoop out/Animals.cs
dotnet run -- --help
dotnet run -- --version
\end{verbatim}

\subsection{Error Handling}

Compile errors throw \cs{XoopCompileException} with a descriptive message.
The CLI prints the message in red to \texttt{stderr} and exits with code
\texttt{1}.

Common errors:
\begin{itemize}[nosep]
  \item Malformed XML (unclosed tags, bad encoding, etc.)
  \item Missing \required\ attributes (\attr{name}, \attr{type}, etc.)
  \item Root element is not \el{XoopProgram}
\end{itemize}

% ═════════════════════════════════════════════════════════════════════════════
\section{Element Quick Reference}

\begin{longtable}{lll}
\toprule
\textbf{Element} & \textbf{Key attributes} & \textbf{C\# output} \\
\midrule
\endfirsthead
\toprule
\textbf{Element} & \textbf{Key attributes} & \textbf{C\# output} \\
\midrule
\endhead
\el{XoopProgram}  & —                                   & (file root) \\
\el{Using}        & \attr{namespace}                    & \cs{using \ldots;} \\
\el{Namespace}    & \attr{name}                         & \cs{namespace \{ \}} \\
\el{Class}        & \attr{name}, \attr{extends}, \ldots & \cs{class / abstract class / static class} \\
\el{Interface}    & \attr{name}, \attr{extends}         & \cs{interface} \\
\el{Enum}         & \attr{name}, \attr{type}            & \cs{enum} \\
\el{Member}       & \attr{name}, \attr{value}           & enum member \\
\el{Field}        & \attr{name}, \attr{type}            & field declaration \\
\el{Property}     & \attr{name}, \attr{type}            & property (auto / backed / custom) \\
\el{Get}          & (body text)                         & \cs{get \{ \}} \\
\el{Set}          & (body text)                         & \cs{set \{ \}} \\
\el{Constructor}  & \attr{access}                       & \cs{ClassName(\ldots) \{ \}} \\
\el{Base}         & \el{Arg} children                  & \cs{: base(\ldots)} \\
\el{This}         & \el{Arg} children                  & \cs{: this(\ldots)} \\
\el{Method}       & \attr{name}, \attr{returnType}      & method declaration \\
\el{Parameter}    & \attr{name}, \attr{type}            & method/ctor parameter \\
\el{Body}         & (body text / CDATA)                 & method body \cs{\{ \}} \\
\el{Arg}          & (text content)                      & argument in \cs{base()}/\cs{this()} \\
\bottomrule
\end{longtable}

% ═════════════════════════════════════════════════════════════════════════════
\section{Design Notes \& Limitations}

\begin{enumerate}
  \item \textbf{Bodies are opaque C\#.} The compiler does not parse or validate
        method bodies — they are emitted verbatim (after whitespace normalisation).
        Type errors will surface at C\# compile time.

  \item \textbf{Generics constraints} (\cs{where T : class}) are not yet
        supported as first-class attributes; write them directly in the body
        if needed, or use a comment to document intent.

  \item \textbf{Attributes / annotations} (C\# \cs{[Attribute]}) are not yet
        supported. A future \el{Annotation} element is planned.

  \item \textbf{Properties with \attr{set=``false''}} emit only a getter.
        A future \attr{init} accessor will be added for init-only setters.

  \item \textbf{Whitespace normalisation} — the emitter strips the minimum
        common leading indent from body text so that naturally indented CDATA
        sections produce correctly indented C\# output.

  \item \textbf{No runtime dependency} — the emitted C\# references only the
        standard .NET base class library. XOOP itself is a compile-time tool only.
\end{enumerate}

\vfill
\begin{center}
  \small\textcolor{xoop-gray}{XOOP Compiler v1.0.0 \quad\textbullet\quad
  Generated \today \quad\textbullet\quad
  Compiles to C\# \textbullet\quad Zero runtime dependency}
\end{center}

\end{document}
